\chapter{Analisis}

\section{Analisis Preeliminar}
El proceso de gestión de nómina en Café y Vivero Verde Pradera se realiza de
forma manual, lo que ha generado errores frecuentes en los cálculos salariales, pérdida
tiempo en tareas repetitivas y dificultades en la emisión de reportes oficiales. El
sistema propuesto busca automatizar estas operaciones, mejorar la precisión del proceso,
optimizar el tiempo y proteger los datos sensibles de los empleados. Además, facilitará
la consulta remota de la información mediante una plataforma web segura y accesible.

\section{Estudio de factibilidad}

\subsection{Factibilidad Operativa}

\begin{itemize}
	\item \textbf{Conocimiento de los desarrolladores:} Los desarrolladores del sistema son estudiantes avanzados de Ingeniería en Sistemas, con experiencia en desarrollo web, bases de datos y metodologías ágiles, lo que garantiza una correcta implementación del proyecto.

	\item \textbf{Habilidad de los usuarios:} El sistema está orientado al personal administrativo de la empresa, quienes ya realizan los procesos manuales y tienen el conocimiento necesario para adaptarse al uso de un sistema digital. Se planea además incluir una etapa de capacitación.

	\item \textbf{Disponibilidad de horarios:} Tanto el equipo de desarrollo como el personal administrativo disponen de horarios establecidos para reuniones de revisión, pruebas e implementación.

	\item \textbf{Mantenimiento del software:} El sistema será documentado técnicamente y se mantendrá en un repositorio GitHub, lo cual facilitará futuras mejoras, corrección de errores y mantenimiento general.
\end{itemize}

\textbf{Conclusión:} Existe una alta viabilidad operativa para implementar el sistema, debido a la disposición de los usuarios, la experiencia del equipo desarrollador y la estructura de apoyo brindada por la empresa.

\subsection{Factibilidad Técnica}

\begin{itemize}
	\item \textbf{Especificación de Hardware y Software:} El sistema funcionará sobre servidores con al menos 4~GB de RAM y 20~GB de almacenamiento. Los usuarios accederán desde navegadores modernos como Chrome o Edge.

	\item \textbf{Infraestructura de redes:} La empresa cuenta con conexión a internet estable, lo cual permite que el sistema funcione correctamente en modalidad web.

	\item \textbf{Herramientas de desarrollo:} Se utilizarán herramientas como Visual Studio Code, GitHub, PostgreSQL, Next.js, Node.js, Prisma ORM, todas gratuitas y ampliamente documentadas.

	\item \textbf{Plataformas de desarrollo:} El sistema será construido como una aplicación web, compatible con sistemas operativos modernos. El backend correrá en Node.js y el frontend en Next.js con TypeScript.
\end{itemize}

\textbf{Conclusión:} El proyecto es técnicamente viable gracias a la disponibilidad de herramientas, infraestructura básica y conocimientos del equipo para implementar la solución sin depender de tecnologías costosas o inaccesibles.

\subsection{Factibilidad Económica}

\begin{itemize}
	\item \textbf{Costos de desarrollo:} Al ser un proyecto académico, no se incurrirá en costos por contratación de programadores o diseñadores.

	\item \textbf{Costos de capital humano:} Los estudiantes desarrolladores participan como parte del trabajo final de graduación. No se requiere remuneración económica directa, lo cual reduce significativamente los costos.

	\item \textbf{Cargos administrativos:} Los recursos administrativos necesarios para colaborar con el proyecto (reuniones, validaciones) serán asumidos por la empresa como parte de su operación normal.

	\item \textbf{Costos totales del proyecto:} El proyecto no requiere inversión económica significativa, ya que se usará software libre y recursos existentes. Se considera un sistema de bajo costo y alto retorno.
\end{itemize}

\textbf{Conclusión:} El proyecto es completamente viable desde el punto de vista económico. Requiere una inversión mínima y promete beneficios importantes como la reducción de errores y ahorro de tiempo en el cálculo de nómina.

\section{Ciclo de vida}

Para el desarrollo del sistema de gestión de nóminas para Café y Vivero Verde Pradera se aplicará un enfoque ágil, utilizando la metodología \textbf{Scrum}, complementada con un modelo de \textbf{Desarrollo Incremental}. El proceso se organizará en ocho sprints consecutivos, cada uno con una duración estimada de dos semanas.

Cada sprint incluye una etapa de revisión y retroalimentación (\textit{Feedback}), lo cual permitirá realizar mejoras iterativas y garantizar que las funcionalidades desarrolladas respondan a las necesidades reales del cliente.

\subsection*{Resumen de los Sprints}

\begin{itemize}
	\item \textbf{Sprint 1: Setup del Proyecto} \\
	      Inicialización del entorno de desarrollo, configuración del repositorio GitHub y definición de la estructura base del sistema.

	\item \textbf{Sprint 2: Módulo de Login} \\
	      Desarrollo del sistema de autenticación de usuarios y control de acceso.

	\item \textbf{Sprint 3: Módulo de Interfaz de Usuario} \\
	      Implementación del diseño general del frontend y componentes base de navegación.

	\item \textbf{Sprint 4: Administración de Eventos Laborales} \\
	      Funcionalidad para registrar y gestionar vacaciones, incapacidades y permisos.

	\item \textbf{Sprint 5: Cálculo de Nóminas} \\
	      Implementación del motor de cálculo automático de salarios y deducciones.

	\item \textbf{Sprint 6: Generación de Reportes} \\
	      Emisión de reportes institucionales y comprobantes de pago para los empleados.

	\item \textbf{Sprint 7: Módulo de Seguridad} \\
	      Validaciones adicionales, cifrado de contraseñas, roles y permisos de usuario.

	\item \textbf{Sprint 8: Pruebas y Entrega} \\
	      Pruebas funcionales, validación con usuarios, documentación final y entrega del sistema.
\end{itemize}

\newpage

\section{Estructura de Desglose del Trabajo (EDT)}

La Estructura de Desglose del Trabajo (EDT) permite organizar el proyecto en fases y
tareas específicas de manera jerárquica. En este caso, se divide en cinco fases principales
que agrupan las actividades necesarias para el desarrollo del sistema de gestión de nóminas.
La siguiente figura representa gráficamente esta estructura, facilitando la planificación y el seguimiento del proyecto.


\begin{figure}[h!]
	\centering
	\begin{tikzpicture}[
			scale=0.75, transform shape,
			font=\small\sffamily,
			box/.style={draw, rounded corners=3pt, text width=3.2cm, align=center, minimum height=1cm, drop shadow, fill=gray!20},
			phase/.style={draw, rounded corners=3pt, text width=3.2cm, align=center, minimum height=1cm, fill=blue!20},
			task/.style={draw, rounded corners=3pt, text width=3.2cm, align=center, minimum height=1cm, fill=orange!20},
			node distance=1cm
		]

		% Fases principales (nivel 2)
		\node[phase] (p1) at (0,0) {1. Dirección del Proyecto};
		\node[phase] (p2) [right=of p1] {2. Recolección y Análisis};
		\node[phase] (p3) [right=of p2] {3. Diseño del Sistema};
		\node[phase] (p4) [right=of p3] {4. Construcción};
		\node[phase] (p5) [right=of p4] {5. Pruebas e Implementación};

		% Nodo raíz centrado respecto a p3
		\node[box] (root) [above=2cm of p3] {Sistema de Gestión de Nóminas};

		% Flechas del nodo raíz
		\foreach \i in {1,...,5} {
				\draw[->, thick] (root.south) -- (p\i.north);
			}
		% Tareas de p1
		\node[task] (p1a) [below=of p1] {Planificación};
		\node[task] (p1b) [below=of p1a] {Reuniones};
		\node[task] (p1c) [below=of p1b] {Administración};
		\draw[->] (p1.south) -- (p1a.north);
		\draw[->] (p1a.south) -- (p1b.north);
		\draw[->] (p1b.south) -- (p1c.north);

		% Tareas de p2
		\node[task] (p2a) [below=of p2] {Requerimientos del sistema};
		\node[task] (p2b) [below=of p2a] {Estudio de factibilidad};
		\node[task] (p2c) [below=of p2b] {Historias de usuario};
		\draw[->] (p2.south) -- (p2a.north);
		\draw[->] (p2a.south) -- (p2b.north);
		\draw[->] (p2b.south) -- (p2c.north);

		% Tareas de p3
		\node[task] (p3a) [below=of p3] {Diseño de base de datos};
		\node[task] (p3b) [below=of p3a] {Frontend y API};
		\node[task] (p3c) [below=of p3b] {Prototipos de interfaz};
		\draw[->] (p3.south) -- (p3a.north);
		\draw[->] (p3a.south) -- (p3b.north);
		\draw[->] (p3b.south) -- (p3c.north);

		% Tareas de p4
		\node[task] (p4a) [below=of p4] {Programación del sistema};
		\node[task] (p4b) [below=of p4a] {Integración de módulos};
		\node[task] (p4c) [below=of p4b] {Documentación técnica};
		\draw[->] (p4.south) -- (p4a.north);
		\draw[->] (p4a.south) -- (p4b.north);
		\draw[->] (p4b.south) -- (p4c.north);

		% Tareas de p5
		\node[task] (p5a) [below=of p5] {Pruebas funcionales};
		\node[task] (p5b) [below=of p5a] {Capacitación de usuarios};
		\node[task] (p5c) [below=of p5b] {Documentación final};
		\draw[->] (p5.south) -- (p5a.north);
		\draw[->] (p5a.south) -- (p5b.north);
		\draw[->] (p5b.south) -- (p5c.north);

	\end{tikzpicture}
	\caption{Estructura de Desglose del Trabajo del Proyecto (EDT)}
\end{figure}

\section{Historias de Usuario}

Las historias de usuario permiten capturar de manera sencilla y centrada en el
usuario las funcionalidades esperadas del sistema de gestión de nóminas para Café
y Vivero Verde Pradera. Estas historias han sido redactadas desde la perspectiva de
los distintos tipos de usuarios del sistema (administradores, colaboradores y gerentes),
con el fin de comprender sus necesidades, priorizar el desarrollo de funcionalidades y
planificar el trabajo en sprints.

Cada historia incluye información clave como el usuario involucrado, la prioridad del
requerimiento, el riesgo asociado, la iteración en que se desarrollará, la estimación
de esfuerzo y una breve descripción del comportamiento esperado del sistema.
Estas historias se utilizarán como base para guiar el diseño, implementación y
pruebas de cada funcionalidad a lo largo del desarrollo del proyecto.

\begin{table}[h]
	\centering
	\caption{Historia de Usuario HU-01}
	\begin{tabular}{|p{5cm}|p{9cm}|}
		\hline
		\textbf{Número:}                          & HU-01                                                                                                                                                                                                                                                                                                       \\ \hline
		\textbf{Nombre:}                          & Registro de usuarios                                                                                                                                                                                                                                                                                        \\ \hline
		\textbf{Usuario:}                         & Administrador                                                                                                                                                                                                                                                                                               \\ \hline
		\textbf{Modificación de Historia Número:} & Nueva                                                                                                                                                                                                                                                                                                       \\ \hline
		\textbf{Prioridad en Negocio:}            & Alta                                                                                                                                                                                                                                                                                                        \\ \hline
		\textbf{Riesgo en Desarrollo:}            & Medio                                                                                                                                                                                                                                                                                                       \\ \hline
		\textbf{Iteración Asignada:}              & Sprint 2: Módulo de Login                                                                                                                                                                                                                                                                                   \\ \hline
		\textbf{Puntos Estimados:}                & 5                                                                                                                                                                                                                                                                                                           \\ \hline
		\textbf{Puntos Reales:}                   & --                                                                                                                                                                                                                                                                                                          \\ \hline
		\textbf{Descripción:}                     & En la administración del sistema se tendrá la opción de registrar usuarios. El administrador podrá ingresar a esta opción, llenar los datos requeridos como nombre, correo y rol, y luego guardar el nuevo usuario. El sistema mostrará una confirmación y almacenará el nuevo usuario en la base de datos. \\ \hline
		\textbf{Observaciones:}                   & El correo debe validarse.                                                                                                                                                                                                                                                                                   \\ \hline
	\end{tabular}
\end{table}


\begin{table}[h]
	\centering
	\caption{Historia de Usuario HU-02}
	\begin{tabular}{|p{5cm}|p{9cm}|}
		\hline
		\textbf{Número:}                          & HU-02                                                                                                                                                                                                                                           \\ \hline
		\textbf{Nombre:}                          & Asignar roles                                                                                                                                                                                                                                   \\ \hline
		\textbf{Usuario:}                         & Administrador                                                                                                                                                                                                                                   \\ \hline
		\textbf{Modificación de Historia Número:} & Nueva                                                                                                                                                                                                                                           \\ \hline
		\textbf{Prioridad en Negocio:}            & Alta                                                                                                                                                                                                                                            \\ \hline
		\textbf{Riesgo en Desarrollo:}            & Bajo                                                                                                                                                                                                                                            \\ \hline
		\textbf{Iteración Asignada:}              & Sprint 3: Módulo de Interfaz de Usuario                                                                                                                                                                                                         \\ \hline
		\textbf{Puntos Estimados:}                & 5                                                                                                                                                                                                                                               \\ \hline
		\textbf{Puntos Reales:}                   & --                                                                                                                                                                                                                                              \\ \hline
		\textbf{Descripción:}                     & El administrador puede asignar un rol a cada usuario desde la administración del sistema. Para ello, selecciona un usuario y el sistema despliega los roles disponibles. Al seleccionar uno, el cambio se guarda y se muestra una confirmación. \\ \hline
		\textbf{Observaciones:}                   & Debe validarse que el usuario tenga solo un rol activo.                                                                                                                                                                                         \\ \hline
	\end{tabular}
\end{table}


\begin{table}[h]
	\centering
	\caption{Historia de Usuario HU-03}
	\begin{tabular}{|p{5cm}|p{9cm}|}
		\hline
		\textbf{Número:}                          & HU-03                                                                                                                                                       \\ \hline
		\textbf{Nombre:}                          & Visualizar historial de pagos                                                                                                                               \\ \hline
		\textbf{Usuario:}                         & Colaborador                                                                                                                                                 \\ \hline
		\textbf{Modificación de Historia Número:} & Nueva                                                                                                                                                       \\ \hline
		\textbf{Prioridad en Negocio:}            & Media                                                                                                                                                       \\ \hline
		\textbf{Riesgo en Desarrollo:}            & Bajo                                                                                                                                                        \\ \hline
		\textbf{Iteración Asignada:}              & Sprint 5: Cálculo de Nóminas                                                                                                                                \\ \hline
		\textbf{Puntos Estimados:}                & 5                                                                                                                                                           \\ \hline
		\textbf{Puntos Reales:}                   & --                                                                                                                                                          \\ \hline
		\textbf{Descripción:}                     & El colaborador accede al historial desde su perfil. Puede ver los montos y fechas de los pagos realizados. El sistema permite ordenar y filtrar por fechas. \\ \hline
		\textbf{Observaciones:}                   & Asegurar que los datos estén protegidos.                                                                                                                    \\ \hline
	\end{tabular}
\end{table}


\begin{table}[h]
	\centering
	\caption{Historia de Usuario HU-04}
	\begin{tabular}{|p{5cm}|p{9cm}|}
		\hline
		\textbf{Número:}                          & HU-04                                                                                                                                                                                   \\ \hline
		\textbf{Nombre:}                          & Generar reporte de nómina                                                                                                                                                               \\ \hline
		\textbf{Usuario:}                         & Gerente                                                                                                                                                                                 \\ \hline
		\textbf{Modificación de Historia Número:} & Nueva                                                                                                                                                                                   \\ \hline
		\textbf{Prioridad en Negocio:}            & Alta                                                                                                                                                                                    \\ \hline
		\textbf{Riesgo en Desarrollo:}            & Alto                                                                                                                                                                                    \\ \hline
		\textbf{Iteración Asignada:}              & Sprint 6: Generación de Reportes                                                                                                                                                        \\ \hline
		\textbf{Puntos Estimados:}                & 5                                                                                                                                                                                       \\ \hline
		\textbf{Puntos Reales:}                   & --                                                                                                                                                                                      \\ \hline
		\textbf{Descripción:}                     & El gerente puede generar un reporte mensual de nómina. Desde el módulo correspondiente, selecciona el mes y el sistema crea un PDF con los detalles de salarios, deducciones y totales. \\ \hline
		\textbf{Observaciones:}                   & Validar acceso solo para gerentes.                                                                                                                                                      \\ \hline
	\end{tabular}
\end{table}

\clearpage


\section{Especificación de Requerimientos de Software (SRS)}
\subsection*{Introducción}
Esta Especificación de Requerimientos de Software (SRS) describe los
requerimientos funcionales y no funcionales del sistema web de gestión de nómina
para la empresa Café y Vivero Verde Pradera. Está dirigida al equipo de desarrollo,
usuarios finales, patrocinadores del proyecto y cualquier persona involucrada en la
validación y verificación del sistema.

\subsubsection*{Propósito}
El sistema gestionará de forma automatizada el cálculo de nómina, eventos laborales,
reportes institucionales y seguridad de datos del personal. Será accesible mediante
navegador web, con controles de acceso para usuarios autorizados. No incluirá
funcionalidades de facturación, contabilidad general o inventario.

\subsubsection*{Alcance}

El sistema gestionará de forma automatizada el cálculo de nómina,
eventos laborales, reportes institucionales y seguridad de datos del
personal. Será accesible mediante navegador web, con controles de acceso
para usuarios autorizados. No incluirá funcionalidades de facturación,
contabilidad general o inventario.

\subsubsection*{Definiciones, siglas y abreviaturas}

\begin{itemize}
	\item \textbf{SRS:} \textit{Software Requirements Specification}.
	\item \textbf{CCSS:} \textit{Caja Costarricense del Seguro Social}.
	\item \textbf{MTSS:} \textit{Ministerio de Trabajo y Seguridad Social}.
	\item \textbf{Nómina:} Proceso de cálculo de salarios y deducciones laborales.
	\item \textbf{Scrum:} Marco de trabajo ágil para gestión de proyectos.
	\item \textbf{VPG\_:} Prefijo usado para nombrar tablas en la base de datos.
\end{itemize}

\subsubsection*{Referencias}

\begin{itemize}
	\item \textbf{TFM del proyecto:} “Sistema de Gestión de Nóminas para Café y Vivero Verde Pradera”. Universidad Nacional, Sede Regional Brunca.

	\item \textbf{Scrum Guide (2020):} Guía oficial de la metodología Scrum por Ken Schwaber y Jeff Sutherland.

	\item \textbf{Normativa de la CCSS:} Requisitos para la generación y entrega de reportes laborales a la Caja Costarricense del Seguro Social.

	\item \textbf{Ministerio de Trabajo y Seguridad Social de Costa Rica:} Lineamientos legales sobre remuneraciones, horarios laborales, licencias y documentación obligatoria relacionada con la gestión de personal.
\end{itemize}


\subsubsection*{Visión general del documento}
\begin{itemize}
	\item \textbf{Sección 2:} Proporciona una descripción general del sistema, incluyendo su perspectiva, funciones principales, usuarios y restricciones de diseño.

	\item \textbf{Sección 3:} Detalla los requerimientos específicos, incluyendo casos de uso del sistema y requerimientos suplementarios.

	\item \textbf{Sección 4:} Describe los requerimientos de documentación que acompañarán al sistema.

\end{itemize}

\subsection*{Descripción General}

Esta sección presenta una visión global del sistema de gestión de nóminas para
Café y Vivero Verde Pradera. Aquí se describen los factores generales que influyen en
el desarrollo del software, tales como el entorno operativo, los usuarios involucrados,
las funciones principales y las restricciones de diseño. Aunque esta sección no entra en
detalle sobre los requerimientos específicos, sienta las bases conceptuales y técnicas que
guían el resto del documento.

\subsubsection*{Perspectiva del producto}

El sistema de gestión de nóminas es una aplicación web que funcionará como un
producto independiente, aunque podría integrarse a futuro con otros sistemas de gestión
empresarial. Su diseño sigue el modelo de n capas , con acceso controlado mediante
credenciales de usuario. El sistema será accesible mediante navegadores modernos desde
cualquier ubicación con conexión a internet.

\subsubsection*{Interfaces de usuario}

\begin{itemize}
	\item Pantalla de inicio de sesión
	\item Panel principal con resumen de información
	\item Módulo de gestión de empleados
	\item Módulo de cálculo y visualización de nómina
	\item Módulo de eventos laborales (vacaciones, incapacidades, permisos)
	\item Módulo de reportes (PDF/Excel)
	\item Módulo de configuración del sistema
	\item Mensajes de validación y error amigables
\end{itemize}

\subsubsection*{Interfaces de Hardware}

\begin{itemize}
	\item El sistema requerirá interacción con impresoras para generar documentos
	      físicos como comprobantes de pago, reportes y formularios.

	\item Se requerirá la integración con un reloj marcador, que permitirá subir
	      automáticamente los registros de entrada y salida del personal al sistema.
	      Estos datos serán utilizados para el cálculo de la nómina, especialmente en
	      el control de jornadas laborales, horas extra e inasistencias.

	\item La integración podrá realizarse mediante archivos de exportación (CSV, Excel)
	      generados por el reloj marcador, o mediante una conexión directa a través de API,
	      dependiendo del modelo del dispositivo.

\end{itemize}


\subsubsection*{Interfaces con software}

\begin{itemize}
	\item \textbf{Navegadores compatibles:} Chrome, Firefox, Microsoft Edge.

	\item \textbf{Base de datos:} Se utilizará una base de datos relacional, con preferencia por PostgreSQL o SQL Server, según disponibilidad y requisitos del entorno de producción.

	\item \textbf{Frameworks y tecnologías principales:}
	      \begin{itemize}
		      \item \textbf{Frontend:} Next.js con TypeScript, para el desarrollo de una aplicación web moderna, escalable y optimizada.
		      \item \textbf{Backend:} Node.js con Express.js, para construir una API REST robusta y eficiente que conecte con la base de datos y gestione la lógica de negocio.
	      \end{itemize}

	\item \textbf{Librerías complementarias:}
	      \begin{itemize}
		      \item Generación de reportes en PDF (por ejemplo, PDFKit o jsPDF).
		      \item Prisma ORM para la gestión de base de datos.
		      \item Librerías de validación, autenticación y control de acceso.
	      \end{itemize}
\end{itemize}


\subsubsection*{Interfaces de comunicación}

\begin{itemize}
	\item HTTP/HTTPS como protocolo de comunicación principal.
	\item APIs REST para el intercambio de datos entre frontend y backend.
	\item Encriptación SSL/TLS para comunicaciones seguras.
\end{itemize}

\subsubsection*{Restricciones de memoria}

El sistema está diseñado para operar eficientemente en servidores con un mínimo de:

\begin{itemize}
	\item Memoria RAM: 4GB
	\item Almacenamiento: 20GB
\end{itemize}

Los navegadores cliente deben tener al menos 1GB de RAM disponible.

\subsubsection*{Requerimientos de adecuación al entorno}

\begin{itemize}
	\item El sistema se adapta a diferentes dispositivos (responsive), permitiendo su uso en computadoras, laptops y tablets.
	\item Idioma predeterminado: Español (Costa Rica)
	\item El servidor deberá estar configurado con zona horaria GMT-6.
\end{itemize}

\subsubsection*{Funciones del producto (Casos de uso de alto nivel)}

\begin{itemize}
	\item Autenticación de usuarios
	\item Gestión de datos de empleados
	\item Cálculo automático de nómina
	\item Registro de eventos laborales (vacaciones, incapacidades, etc.)
	\item Generación de reportes oficiales
	\item Visualización de comprobantes de pago
	\item Control de accesos y seguridad
	\item Capacitación y ayuda al usuario
\end{itemize}

\subsubsection*{Características de los usuarios}

\begin{itemize}
	\item \textbf{Administrador:} Usuario principal del sistema, con experiencia en gestión administrativa. Es responsable de registrar y mantener los datos del personal, roles, jornadas, eventos laborales, generar la nómina, administrar usuarios y emitir reportes oficiales. Tiene acceso total a todos los módulos del sistema, incluyendo la configuración de seguridad y la bitácora.

	\item \textbf{Colaborador (Empleado):} Usuario con acceso restringido. Puede consultar su información personal, visualizar y descargar su historial de pagos, y recibir los comprobantes por correo. No tiene acceso a la administración de nómina, ni puede modificar sus datos directamente.

	\item \textbf{Gerente:} Usuario con acceso a los reportes consolidados de nómina. Puede generar documentación institucional (para CCSS, MTSS, contabilidad), consultar datos agregados y validar información histórica. No tiene permisos para editar nómina ni datos del personal.

	\item \textbf{Contabilidad (opcional):} Usuario con acceso exclusivo a los módulos de exportación de datos contables y reportes oficiales. Puede descargar formatos estandarizados y generar documentación legal.
\end{itemize}

\subsubsection*{Restricciones de diseño}

\begin{itemize}
	\item \textbf{Lenguaje backend:} TypeScript utilizando Node.js y Express.js como framework principal para el desarrollo de la API.
	\item \textbf{Lenguaje frontend:} Next.js con TypeScript, que ofrece funcionalidades modernas como renderizado híbrido (SSR/SSG) y excelente integración con React y Tailwind CSS para el diseño visual.
	\item \textbf{Base de datos:} PostgreSQL como base preferente por su robustez y escalabilidad, aunque puede considerarse SQL Server según la infraestructura del cliente.
	\item \textbf{ORM:} Prisma será utilizado como herramienta de mapeo objeto-relacional para facilitar la interacción entre la base de datos y el backend.
	\item \textbf{Arquitectura:} Modelo cliente-servidor basado en APIs REST, con separación clara entre frontend (interfaz de usuario) y backend (lógica de negocio y persistencia).
	\item \textbf{Control de versiones:} Git, utilizando GitHub como plataforma central para la gestión del código fuente, seguimiento de cambios y colaboración mediante Pull Requests.
	\item \textbf{Estandarización en base de datos:} Se utilizará un esquema consistente de nomenclatura, donde las tablas comienzan con el prefijo VPG\_ seguido del nombre en mayúsculas, y las columnas en minúsculas con guion bajo.
\end{itemize}

\subsubsection*{Restricción de diseño 1}

\textbf{Uso obligatorio de Next.js con TypeScript como framework de desarrollo frontend.}
El sistema debe ser desarrollado utilizando el framework Next.js con TypeScript como lenguaje principal, por su capacidad de renderizado híbrido, manejo eficiente del enrutamiento y escalabilidad. Esta decisión es clave para mantener un entorno moderno, seguro y compatible con estándares actuales de desarrollo web.

\subsubsection*{Restricción de diseño 2}

\textbf{Backend implementado en Express.js con Node.js y TypeScript.}
La lógica de negocio del sistema debe ser implementada mediante una API RESTful utilizando Express.js sobre Node.js, con escritura en TypeScript para mejorar la robustez del código. Esta tecnología ha sido elegida por su ligereza, flexibilidad, comunidad activa y compatibilidad con herramientas modernas de desarrollo.

\subsubsection*{Restricción de diseño 3}

\textbf{Uso de PostgreSQL como base de datos principal.}
La base de datos del sistema será implementada en PostgreSQL, salvo que se requiera específicamente el uso de SQL Server por condiciones de infraestructura del cliente. PostgreSQL fue seleccionado por su alto rendimiento, seguridad, soporte para transacciones ACID y amplia compatibilidad con herramientas ORM modernas.

\subsubsection*{Restricción de diseño 4}

\textbf{ORM Prisma para gestión de la base de datos.}
Se utilizará Prisma como herramienta de mapeo objeto-relacional (ORM) para interactuar con la base de datos. Prisma proporciona tipado fuerte, migraciones estructuradas y una experiencia de desarrollo alineada con TypeScript, lo que facilita la validación y mantenimiento del sistema.

\subsubsection*{Restricción de diseño 5}

\textbf{Modelo de arquitectura cliente-servidor basado en APIs REST.}
El sistema adoptará una arquitectura cliente-servidor, donde la interfaz de usuario (frontend) se comunica con el backend a través de una API REST. Esta separación garantiza una mejor escalabilidad, facilita el mantenimiento del código y permite una posible integración futura con otros sistemas.

\subsubsection*{Restricción de diseño 6}

\textbf{Control de versiones obligatorio con Git y repositorio en GitHub.}
Todo el código fuente del sistema será gestionado mediante Git, con uso obligatorio del repositorio GitHub. Se seguirán buenas prácticas como el uso de ramas (main, develop, feature/*), mensajes de commit con prefijos ([FEATURE], [FIXED], [BUG], etc.) y revisiones por medio de pull requests para evitar conflictos y mantener calidad de código.

\subsubsection*{Restricción de diseño 7}

\textbf{Nomenclatura de base de datos estandarizada.}
Las tablas en la base de datos deberán seguir el prefijo VPG\_ en mayúsculas, mientras que los campos se escribirán en minúsculas con guion bajo (snake\_case). Las claves foráneas seguirán el formato VPG\_TABLA1\_TABLA2\_FK. Esto garantiza consistencia, claridad y facilita la lectura del modelo de datos.

\subsubsection*{Supuestos y dependencias}

\begin{itemize}
	\item Se asume que los usuarios tienen conocimientos básicos en el uso de sistemas web.
	\item Se espera que la empresa mantenga infraestructura básica (conectividad, energía eléctrica).
	\item Cualquier cambio en la legislación laboral podría requerir ajustes al sistema.
\end{itemize}


\subsection*{Requerimientos del Sistema}

Esta sección detalla los requerimientos funcionales y no funcionales del sistema
de gestión de nóminas para Café y Vivero Verde Pradera. Los requerimientos se han
organizado utilizando el modelo de casos de uso para representar las funcionalidades
principales, así como especificaciones suplementarias que cubren aspectos como rendimiento,
seguridad, usabilidad y documentación.

La información aquí contenida proporciona el nivel de detalle necesario para que el
equipo de desarrollo diseñe e implemente el sistema, y para que los encargados de calidad
validen su cumplimiento mediante pruebas.

\subsubsection*{Modelo de Caso de Uso de Negocio}

El modelo de negocio se centra en optimizar la gestión de nóminas del personal de la empresa Café y Vivero Verde Pradera. A continuación, se enumeran los principales procesos, definidos con base en las historias de usuario priorizadas:

\begin{itemize}
	\item Ingreso al sistema (autenticación de usuarios)
	\item Registro y administración de empleados
	\item Asignación y gestión de roles de usuario
	\item Registro de jornada laboral (llegada y salida)
	\item Registro y administración de eventos laborales (vacaciones, incapacidades, permisos)
	\item Cálculo automático de salarios y deducciones
	\item Generación de reportes oficiales (CCSS, MTSS, contabilidad)
	\item Emisión de comprobantes de pago a colaboradores
	\item Configuración de parámetros del sistema (pendiente de historia funcional)
\end{itemize}

Estos procesos serán representados como Casos de Uso en el sistema de información.

\subsubsection*{Modelo de Caso de Uso de Sistema}

\begin{table}[h]
	\centering
	\caption{Modelo de Casos de Uso del Sistema}
	\begin{tabular}{|c|p{6cm}|p{5cm}|}
		\hline
		\textbf{Código} & \textbf{Nombre de Caso de uso}    & \textbf{Actor Principal} \\ \hline
		CU01            & Iniciar Sesión                    & Usuario autorizado       \\ \hline
		CU02            & Registrar empleado                & Administrador            \\ \hline
		CU03            & Editar información de empleado    & Administrador            \\ \hline
		CU04            & Registrar evento laboral          & Administrador            \\ \hline
		CU05            & Calcular nómina                   & Administrador            \\ \hline
		CU06            & Generar reporte para la CCSS      & Administrador            \\ \hline
		CU07            & Consultar comprobante de pago     & Administrador            \\ \hline
		CU08            & Configurar parámetros del sistema & Administrador            \\ \hline
	\end{tabular}
\end{table}

\clearpage

\subsubsection*{Especificación de Caso de Uso: CU01 - Iniciar sesión}

\begin{table}[h]
	\centering
	\caption{Especificación de Caso de Uso CU01 - Iniciar sesión}
	\begin{tabular}{|p{5cm}|p{9cm}|}
		\hline
		\textbf{Especificación de Caso de Uso} & CU01 - Iniciar sesión \\ \hline
		\textbf{Versión}                       & 1.0                   \\ \hline
		\textbf{CU de Negocio}                 & Iniciar sesión        \\ \hline
		\textbf{Fecha}                         & 10/04/2025            \\ \hline
		\textbf{Código del Caso de Uso}        & CU01                  \\ \hline
	\end{tabular}
\end{table}

\subsubsection*{Descripción}
Este caso de uso permite a los usuarios autorizados acceder al sistema mediante autenticación.

\subsubsection*{Precondiciones}
\begin{itemize}
	\item El usuario debe estar registrado previamente en el sistema.
	\item El navegador debe tener conexión activa a internet.
	\item El sistema debe estar disponible.
\end{itemize}

\subsubsection*{Poscondiciones}
\begin{itemize}
	\item El usuario accede al sistema con los permisos correspondientes.
	\item Se registra la hora de acceso en el historial.
\end{itemize}

\subsubsection*{Referencias}
CU02

\clearpage

\subsubsection*{Flujo de Eventos}
\begin{table}[h]
	\centering
	\caption{Flujo principal de eventos}
	\begin{tabular}{|p{7cm}|p{7cm}|}
		\hline
		\textbf{Acción del actor}                     & \textbf{Respuesta del sistema}                              \\ \hline
		El usuario abre la página de inicio de sesión & El sistema muestra el formulario de autenticación           \\ \hline
		El usuario ingresa sus credenciales           & El sistema valida las credenciales                          \\ \hline
		El usuario hace clic en 'Iniciar sesión'      & El sistema concede acceso si las credenciales son correctas \\ \hline
	\end{tabular}
\end{table}

\subsubsection*{Flujos alternos}
\begin{table}[h]
	\centering
	\caption{Flujos alternos}
	\begin{tabular}{|p{7cm}|p{7cm}|}
		\hline
		\textbf{Acción del actor}                     & \textbf{Respuesta del sistema}                                       \\ \hline
		El usuario deja campos vacíos                 & El sistema muestra un mensaje de error indicando campos obligatorios \\ \hline
		El usuario introduce credenciales incorrectas & El sistema muestra 'Credenciales inválidas'                          \\ \hline
	\end{tabular}
\end{table}

\subsubsection*{Excepciones}
\begin{table}[h]
	\centering
	\caption{Excepciones del caso de uso CU01}
	\begin{tabular}{|p{4cm}|p{6cm}|p{6cm}|}
		\hline
		\textbf{Código Excepción} & \textbf{Evento}      & \textbf{Excepción}                        \\ \hline
		E001                      & Servidor no responde & Mostrar mensaje: 'Servicio no disponible' \\ \hline
		E002                      & Usuario bloqueado    & Mostrar advertencia y denegar acceso      \\ \hline
	\end{tabular}
\end{table}

\subsubsection*{Estado}
En definición

\subsubsection*{Especificación de Caso de Uso: CU02 - Registrar empleado}

\begin{table}[h]
	\centering
	\caption{Especificación de Caso de Uso CU02 - Registrar empleado}
	\begin{tabular}{|p{5cm}|p{9cm}|}
		\hline
		\textbf{Especificación de Caso de Uso} & CU02 - Registrar empleado \\ \hline
		\textbf{Versión}                       & 1.0                       \\ \hline
		\textbf{CU de Negocio}                 & Registro de empleados     \\ \hline
		\textbf{Fecha}                         & 10/04/2025                \\ \hline
		\textbf{Código del Caso de Uso}        & CU02                      \\ \hline
	\end{tabular}
\end{table}

\subsubsection*{Descripción}
Este caso de uso permite al administrador registrar nuevos empleados en el sistema, ingresando datos personales, laborales y de contacto. Esta información se usará en procesos de nómina, reportes y control administrativo.

\subsubsection*{Precondiciones}
\begin{itemize}
	\item El usuario debe haber iniciado sesión como administrador.
	\item El sistema debe estar conectado a la base de datos.
	\item El formulario de registro debe estar disponible.
\end{itemize}

\subsubsection*{Poscondiciones}
\begin{itemize}
	\item El nuevo empleado queda registrado correctamente en la base de datos.
	\item La información se encuentra disponible para ser consultada y utilizada por otros módulos.
\end{itemize}

\subsubsection*{Referencias}
CU01, CU03, CU05

\clearpage

\subsubsection*{Flujo de Eventos}
\begin{table}[h]
	\centering
	\caption{Flujo principal de eventos}
	\begin{tabular}{|p{7cm}|p{7cm}|}
		\hline
		\textbf{Acción del actor}                        & \textbf{Respuesta del sistema}                       \\ \hline
		El administrador selecciona “Registrar empleado” & El sistema muestra el formulario de registro         \\ \hline
		El administrador llena los campos requeridos     & El sistema valida los datos ingresados               \\ \hline
		El administrador hace clic en “Guardar”          & El sistema almacena los datos y confirma el registro \\ \hline
	\end{tabular}
\end{table}

\subsubsection*{Flujos alternos}
\begin{table}[h]
	\centering
	\caption{Flujos alternos}
	\begin{tabular}{|p{7cm}|p{7cm}|}
		\hline
		\textbf{Acción del actor}                               & \textbf{Respuesta del sistema}                                     \\ \hline
		El administrador deja campos obligatorios sin completar & El sistema muestra mensaje de error y resalta los campos faltantes \\ \hline
		El administrador ingresa una cédula ya registrada       & El sistema muestra advertencia: "Empleado ya registrado"           \\ \hline
	\end{tabular}
\end{table}

\subsubsection*{Excepciones}
\begin{table}[h]
	\centering
	\caption{Excepciones del caso de uso CU02}
	\begin{tabular}{|p{3cm}|p{5cm}|p{5cm}|}
		\hline
		\textbf{Código Excepción} & \textbf{Evento}                        & \textbf{Excepción}                                       \\ \hline
		E001                      & Fallo de conexión con la base de datos & Mostrar mensaje: 'No se pudo guardar, intente más tarde' \\ \hline
		E002                      & Formato incorrecto de algún dato       & Mostrar advertencia y evitar el guardado                 \\ \hline
	\end{tabular}
\end{table}

\newpage

\subsubsection*{Estado}
En definición

\subsubsection*{Especificación de Caso de Uso: CU03 - Editar información de empleado}

\begin{table}[h]
	\centering
	\caption{Especificación de Caso de Uso CU03 - Editar información de empleado}
	\begin{tabular}{|p{5cm}|p{9cm}|}
		\hline
		\textbf{Especificación de Caso de Uso} & CU03 - Editar información de empleado \\ \hline
		\textbf{Versión}                       & 1.0                                   \\ \hline
		\textbf{CU de Negocio}                 & Registro de empleados                 \\ \hline
		\textbf{Fecha}                         & 10/04/2025                            \\ \hline
		\textbf{Código del Caso de Uso}        & CU03                                  \\ \hline
	\end{tabular}
\end{table}

\subsubsection*{Descripción}
Este caso de uso permite al administrador modificar los datos personales o laborales de un empleado previamente registrado. Las modificaciones se reflejan en los procesos de cálculo de nómina, reportes y control administrativo.

\subsubsection*{Precondiciones}
\begin{itemize}
	\item El usuario debe haber iniciado sesión como administrador.
	\item Debe existir al menos un empleado registrado en el sistema.
	\item El sistema debe estar conectado a la base de datos.
\end{itemize}

\subsubsection*{Poscondiciones}
\begin{itemize}
	\item Los datos del empleado quedan actualizados en el sistema.
	\item Las nuevas configuraciones impactan en cálculos y reportes futuros.
\end{itemize}

\subsubsection*{Referencias}
CU01, CU02, CU05

\clearpage

\subsubsection*{Flujo de Eventos}
\begin{table}[h]
	\centering
	\caption{Flujo principal de eventos}
	\begin{tabular}{|p{7cm}|p{7cm}|}
		\hline
		\textbf{Acción del actor}                       & \textbf{Respuesta del sistema}                            \\ \hline
		El administrador selecciona un empleado         & El sistema carga la información actual del empleado       \\ \hline
		El administrador modifica los campos deseados   & El sistema valida los datos modificados                   \\ \hline
		El administrador hace clic en “Guardar cambios” & El sistema actualiza los datos y muestra mensaje de éxito \\ \hline
	\end{tabular}
\end{table}

\subsubsection*{Flujos alternos}
\begin{table}[h]
	\centering
	\caption{Flujos alternos}
	\begin{tabular}{|p{7cm}|p{7cm}|}
		\hline
		\textbf{Acción del actor}                          & \textbf{Respuesta del sistema}                                  \\ \hline
		El administrador intenta guardar sin hacer cambios & El sistema muestra mensaje: "No se detectaron cambios"          \\ \hline
		El administrador deja campos obligatorios vacíos   & El sistema muestra advertencia y resalta los campos incompletos \\ \hline
	\end{tabular}
\end{table}

\subsubsection*{Excepciones}
\begin{table}[h]
	\centering
	\caption{Excepciones del caso de uso CU03}
	\begin{tabular}{|p{3cm}|p{5cm}|p{5cm}|}
		\hline
		\textbf{Código Excepción} & \textbf{Evento}                         & \textbf{Excepción}                                             \\ \hline
		E001                      & Fallo en la conexión a la base de datos & Mostrar mensaje: 'Error al guardar cambios, intente más tarde' \\ \hline
		E002                      & Formato inválido de datos               & Mostrar advertencia con campo específico                       \\ \hline
	\end{tabular}
\end{table}

\newpage

\subsubsection*{Estado}
En definición

\subsubsection*{Especificación de Caso de Uso: CU04 - Registrar evento laboral}

\begin{table}[h]
	\centering
	\caption{Especificación de Caso de Uso CU04 - Registrar evento laboral}
	\begin{tabular}{|p{5cm}|p{9cm}|}
		\hline
		\textbf{Especificación de Caso de Uso} & CU04 - Registrar evento laboral \\ \hline
		\textbf{Versión}                       & 1.0                             \\ \hline
		\textbf{CU de Negocio}                 & Registro de eventos laborales   \\ \hline
		\textbf{Fecha}                         & 10/04/2025                      \\ \hline
		\textbf{Código del Caso de Uso}        & CU04                            \\ \hline
	\end{tabular}
\end{table}

\subsubsection*{Descripción}
Este caso de uso permite al administrador registrar eventos laborales tales como vacaciones, incapacidades, permisos o ausencias de los empleados. Estos eventos afectan directamente el cálculo de la nómina.

\subsubsection*{Precondiciones}
\begin{itemize}
	\item El usuario debe haber iniciado sesión como administrador.
	\item Debe existir al menos un empleado registrado.
	\item El formulario de eventos laborales debe estar disponible.
\end{itemize}

\subsubsection*{Poscondiciones}
\begin{itemize}
	\item El evento queda registrado y asociado al empleado correspondiente.
	\item Los eventos se reflejan en el cálculo de la nómina.
\end{itemize}

\subsubsection*{Referencias}
CU01, CU02, CU05

\clearpage

\subsubsection*{Flujo de Eventos}
\begin{table}[h]
	\centering
	\caption{Flujo principal de eventos}
	\begin{tabular}{|p{7cm}|p{7cm}|}
		\hline
		\textbf{Acción del actor}                          & \textbf{Respuesta del sistema}                                \\ \hline
		El administrador selecciona un empleado            & El sistema carga el historial de eventos del empleado         \\ \hline
		El administrador escoge el tipo de evento y fechas & El sistema habilita los campos correspondientes               \\ \hline
		El administrador hace clic en “Registrar”          & El sistema guarda el evento y muestra mensaje de confirmación \\ \hline
	\end{tabular}
\end{table}

\subsubsection*{Flujos alternos}
\begin{table}[h]
	\centering
	\caption{Flujos alternos}
	\begin{tabular}{|p{7cm}|p{7cm}|}
		\hline
		\textbf{Acción del actor}                                         & \textbf{Respuesta del sistema}                                                     \\ \hline
		El administrador ingresa un rango de fechas inválido              & El sistema muestra advertencia: "Fecha de fin no puede ser menor que la de inicio" \\ \hline
		El administrador registra un evento duplicado en el mismo periodo & El sistema muestra mensaje: "Ya existe un evento en este periodo"                  \\ \hline
	\end{tabular}
\end{table}

\subsubsection*{Excepciones}
\begin{table}[h]
	\centering
	\caption{Excepciones del caso de uso CU04}
	\begin{tabular}{|p{3cm}|p{5cm}|p{5cm}|}
		\hline
		\textbf{Código Excepción} & \textbf{Evento}                   & \textbf{Excepción}                                               \\ \hline
		E001                      & Error al guardar en base de datos & Mostrar mensaje: 'No se pudo registrar el evento, intente luego' \\ \hline
		E002                      & Formato incorrecto de fecha       & Mostrar advertencia: 'Formato de fecha inválido'                 \\ \hline
	\end{tabular}
\end{table}

\newpage

\subsubsection*{Estado}
En definición

\subsection*{Especificación de Caso de Uso: CU05 - Calcular Nómina}

\begin{table}[h]
	\centering
	\caption{Especificación de Caso de Uso CU05 - Calcular nómina}
	\begin{tabular}{|p{5cm}|p{9cm}|}
		\hline
		\textbf{Especificación de Caso de Uso} & CU05 - Calcular nómina \\ \hline
		\textbf{Versión}                       & 1.0                    \\ \hline
		\textbf{CU de Negocio}                 & Cálculo de Salarios    \\ \hline
		\textbf{Fecha}                         & 08/04/2025             \\ \hline
		\textbf{Código del Caso de Uso}        & CU05                   \\ \hline
	\end{tabular}
\end{table}

\subsubsection*{Descripción}
Este caso de uso permite al administrador calcular la nómina mensual de todos los empleados, tomando en cuenta los eventos laborales registrados, salario base, deducciones obligatorias y horas trabajadas.

\subsubsection*{Precondiciones}
\begin{itemize}
	\item El usuario debe haber iniciado sesión.
	\item Deben existir empleados registrados en el sistema.
	\item Los eventos laborales deben estar actualizados.
\end{itemize}

\subsubsection*{Poscondiciones}
\begin{itemize}
	\item La nómina mensual queda registrada.
	\item Se generan comprobantes de pago individuales por empleado.
	\item Se habilita la generación de reportes para instituciones oficiales.
\end{itemize}

\subsubsection*{Referencias}
CU02, CU03, CU04

\clearpage

\subsubsection*{Flujo de Eventos}

\begin{table}[h]
	\centering
	\caption{Flujo principal de eventos}
	\begin{tabular}{|p{7cm}|p{7cm}|}
		\hline
		\textbf{Acción del actor}                            & \textbf{Respuesta del sistema}                              \\ \hline
		El usuario selecciona el mes y año a calcular        & El sistema carga los datos de empleados y eventos laborales \\ \hline
		El usuario confirma el cálculo                       & El sistema realiza el cálculo automático de salarios        \\ \hline
		El sistema muestra un resumen con los montos totales & Se habilita la opción de guardar y emitir comprobantes      \\ \hline
	\end{tabular}
\end{table}

\subsubsection*{Flujos alternos}

\begin{table}[h]
	\centering
	\caption{Flujos alternos}
	\begin{tabular}{|p{7cm}|p{7cm}|}
		\hline
		\textbf{Acción del actor}                                                   & \textbf{Respuesta del sistema}                                                                                                             \\ \hline
		El usuario selecciona un mes sin registros de eventos laborales             & El sistema muestra un aviso: "No se encontraron eventos laborales para este periodo" y permite continuar el cálculo solo con salarios base \\ \hline
		El usuario intenta calcular la nómina sin haber registrado nuevos empleados & El sistema muestra un mensaje de advertencia: "No hay empleados registrados para calcular nómina" y bloquea el proceso                     \\ \hline
		El usuario modifica manualmente un evento laboral antes del cálculo         & El sistema actualiza automáticamente los datos antes de procesar la nómina                                                                 \\ \hline
	\end{tabular}
\end{table}

\clearpage

\subsubsection*{Excepciones}

\begin{table}[h]
	\centering
	\caption{Excepciones del caso de uso CU05}
	\begin{tabular}{|p{4cm}|p{6cm}|p{6cm}|}
		\hline
		\textbf{Código Excepción} & \textbf{Evento}              & \textbf{Excepción}                     \\ \hline
		E001                      & No hay empleados activos     & Mostrar mensaje de error               \\ \hline
		E002                      & Datos incompletos en eventos & Mostrar advertencia y bloquear cálculo \\ \hline
	\end{tabular}
\end{table}

\subsubsection*{Estado}
En definición

\subsubsection*{Especificación de Caso de Uso: CU06 - Generar reporte para la CCSS}

\begin{table}[h]
	\centering
	\caption{Especificación de Caso de Uso CU06 - Generar reporte para la CCSS}
	\begin{tabular}{|p{5cm}|p{9cm}|}
		\hline
		\textbf{Especificación de Caso de Uso} & CU06 - Generar reporte para la CCSS \\ \hline
		\textbf{Versión}                       & 1.0                                 \\ \hline
		\textbf{CU de Negocio}                 & Generación de reportes oficiales    \\ \hline
		\textbf{Fecha}                         & 10/04/2025                          \\ \hline
		\textbf{Código del Caso de Uso}        & CU06                                \\ \hline
	\end{tabular}
\end{table}

\subsubsection*{Descripción}
Este caso de uso permite al administrador generar el reporte oficial de la nómina de un mes determinado, en el formato requerido por la Caja Costarricense del Seguro Social (CCSS). El reporte puede exportarse en PDF o Excel.

\subsubsection*{Precondiciones}
\begin{itemize}
	\item El usuario debe haber iniciado sesión como administrador.
	\item Debe existir una nómina calculada para el mes seleccionado.
	\item Los datos de empleados y deducciones deben estar actualizados.
\end{itemize}

\subsubsection*{Poscondiciones}
\begin{itemize}
	\item Se genera un archivo descargable con el reporte oficial.
	\item El evento queda registrado en la bitácora del sistema.
\end{itemize}

\subsubsection*{Referencias}
CU01, CU05, RS01

\subsubsection*{Flujo de Eventos}
\begin{table}[h]
	\centering
	\caption{Flujo principal de eventos}
	\begin{tabular}{|p{7cm}|p{7cm}|}
		\hline
		\textbf{Acción del actor}                          & \textbf{Respuesta del sistema}                                 \\ \hline
		El administrador selecciona “Generar reporte CCSS” & El sistema muestra el formulario de selección de mes y formato \\ \hline
		El administrador selecciona el mes y formato       & El sistema valida que exista nómina registrada para ese mes    \\ \hline
		El administrador hace clic en “Generar”            & El sistema genera el archivo y lo presenta para descarga       \\ \hline
	\end{tabular}
\end{table}

\subsubsection*{Flujos alternos}
\begin{table}[h]
	\centering
	\caption{Flujos alternos}
	\begin{tabular}{|p{7cm}|p{7cm}|}
		\hline
		\textbf{Acción del actor}                               & \textbf{Respuesta del sistema}                                             \\ \hline
		El administrador selecciona un mes sin nómina calculada & El sistema muestra advertencia: "No hay datos de nómina para este periodo" \\ \hline
		El administrador selecciona un formato no soportado     & El sistema muestra error: "Formato de exportación no disponible"           \\ \hline
	\end{tabular}
\end{table}

\newpage

\subsubsection*{Excepciones}
\begin{table}[h]
	\centering
	\caption{Excepciones del caso de uso CU06}
	\begin{tabular}{|p{3cm}|p{5cm}|p{5cm}|}
		\hline
		\textbf{Código Excepción} & \textbf{Evento}                    & \textbf{Excepción}                                                 \\ \hline
		E001                      & Error al generar el archivo        & Mostrar mensaje: 'No se pudo generar el archivo, intente de nuevo' \\ \hline
		E002                      & Plantilla de reporte no encontrada & Mostrar mensaje: 'Plantilla de reporte dañada o ausente'           \\ \hline
	\end{tabular}
\end{table}

\subsubsection*{Estado}
En definición


\subsubsection*{Especificación de Caso de Uso: CU07 - Consultar comprobante de pago}

\begin{table}[h]
	\centering
	\caption{Especificación de Caso de Uso CU07 - Consultar comprobante de pago}
	\begin{tabular}{|p{5cm}|p{9cm}|}
		\hline
		\textbf{Especificación de Caso de Uso} & CU07 - Consultar comprobante de pago \\ \hline
		\textbf{Versión}                       & 1.0                                  \\ \hline
		\textbf{CU de Negocio}                 & Emisión de comprobantes de pago      \\ \hline
		\textbf{Fecha}                         & 10/04/2025                           \\ \hline
		\textbf{Código del Caso de Uso}        & CU07                                 \\ \hline
	\end{tabular}
\end{table}

\subsubsection*{Descripción}
Este caso de uso permite tanto al administrador como al empleado (con acceso restringido) consultar los comprobantes de pago generados a partir de la nómina. El sistema permite visualizar e imprimir los comprobantes asociados a cada periodo.

\subsubsection*{Precondiciones}
\begin{itemize}
	\item El usuario debe haber iniciado sesión.
	\item Debe existir al menos un comprobante de pago generado.
	\item El sistema debe tener acceso a los archivos de comprobantes.
\end{itemize}

\subsubsection*{Poscondiciones}
\begin{itemize}
	\item El comprobante es visualizado correctamente por el usuario.
	\item El usuario puede descargar el comprobante en formato PDF.
\end{itemize}

\subsubsection*{Referencias}
CU01, CU05, CU06

\subsubsection*{Flujo de Eventos}
\begin{table}[h]
	\centering
	\caption{Flujo principal de eventos}
	\begin{tabular}{|p{7cm}|p{7cm}|}
		\hline
		\textbf{Acción del actor}                   & \textbf{Respuesta del sistema}                               \\ \hline
		El usuario accede al módulo “Comprobantes”  & El sistema muestra el listado de comprobantes disponibles    \\ \hline
		El usuario selecciona un mes                & El sistema carga el comprobante de pago correspondiente      \\ \hline
		El usuario hace clic en “Ver” o “Descargar” & El sistema presenta el comprobante en pantalla o lo descarga \\ \hline
	\end{tabular}
\end{table}

\subsubsection*{Flujos alternos}
\begin{table}[h]
	\centering
	\caption{Flujos alternos}
	\begin{tabular}{|p{7cm}|p{7cm}|}
		\hline
		\textbf{Acción del actor}                                    & \textbf{Respuesta del sistema}                                     \\ \hline
		El usuario selecciona un mes sin comprobante disponible      & El sistema muestra mensaje: "No hay comprobante para este periodo" \\ \hline
		El usuario intenta abrir un comprobante con error de formato & El sistema muestra mensaje: "Error al cargar el comprobante"       \\ \hline
	\end{tabular}
\end{table}

\subsubsection*{Excepciones}
\begin{table}[h]
	\centering
	\caption{Excepciones del caso de uso CU07}
	\begin{tabular}{|p{3cm}|p{5cm}|p{5cm}|}
		\hline
		\textbf{Código Excepción} & \textbf{Evento}       & \textbf{Excepción}                                           \\ \hline
		E001                      & Archivo no encontrado & Mostrar mensaje: 'Comprobante no disponible o fue eliminado' \\ \hline
		E002                      & Comprobante inválido  & Mostrar advertencia: 'El archivo no se puede visualizar'     \\ \hline
	\end{tabular}
\end{table}

\subsubsection*{Estado}
En definición

\subsubsection*{Especificación de Caso de Uso: CU08 - Configurar parámetros del sistema}

\begin{table}[h]
	\centering
	\caption{Especificación de Caso de Uso CU08 - Configurar parámetros del sistema}
	\begin{tabular}{|p{5cm}|p{9cm}|}
		\hline
		\textbf{Especificación de Caso de Uso} & CU08 - Configurar parámetros del sistema \\ \hline
		\textbf{Versión}                       & 1.0                                      \\ \hline
		\textbf{CU de Negocio}                 & Configuración del sistema                \\ \hline
		\textbf{Fecha}                         & 10/04/2025                               \\ \hline
		\textbf{Código del Caso de Uso}        & CU08                                     \\ \hline
	\end{tabular}
\end{table}

\subsubsection*{Descripción}
Este caso de uso permite al administrador modificar los parámetros globales del sistema, como tasas de deducción, salario mínimo, fechas de corte y opciones visuales. Estos parámetros afectan el comportamiento del sistema en diferentes módulos, incluyendo cálculo de nómina y reportes.

\subsubsection*{Precondiciones}
\begin{itemize}
	\item El usuario debe haber iniciado sesión como administrador.
	\item El sistema debe tener acceso a la base de datos de configuración.
	\item El formulario de configuración debe estar activo.
\end{itemize}

\subsubsection*{Poscondiciones}
\begin{itemize}
	\item Los parámetros quedan actualizados y se aplican al funcionamiento general del sistema.
	\item Las modificaciones se registran en la bitácora del sistema.
\end{itemize}

\subsubsection*{Referencias}
CU01, CU05, CU06

\subsubsection*{Flujo de Eventos}
\begin{table}[h]
	\centering
	\caption{Flujo principal de eventos}
	\begin{tabular}{|p{7cm}|p{7cm}|}
		\hline
		\textbf{Acción del actor}                          & \textbf{Respuesta del sistema}                                 \\ \hline
		El administrador accede al módulo de configuración & El sistema muestra los parámetros actuales del sistema         \\ \hline
		El administrador modifica uno o varios valores     & El sistema valida los nuevos datos ingresados                  \\ \hline
		El administrador hace clic en “Guardar”            & El sistema actualiza los parámetros y muestra mensaje de éxito \\ \hline
	\end{tabular}
\end{table}

\subsubsection*{Flujos alternos}
\begin{table}[h]
	\centering
	\caption{Flujos alternos}
	\begin{tabular}{|p{7cm}|p{7cm}|}
		\hline
		\textbf{Acción del actor}                          & \textbf{Respuesta del sistema}                                     \\ \hline
		El administrador intenta guardar sin hacer cambios & El sistema muestra mensaje: "No se detectaron modificaciones"      \\ \hline
		El administrador ingresa un valor fuera del rango  & El sistema muestra advertencia y bloquea el guardado del parámetro \\ \hline
	\end{tabular}
\end{table}

\subsubsection*{Excepciones}
\begin{table}[h]
	\centering
	\caption{Excepciones del caso de uso CU08}
	\begin{tabular}{|p{3cm}|p{5cm}|p{5cm}|}
		\hline
		\textbf{Código Excepción} & \textbf{Evento}                & \textbf{Excepción}                                          \\ \hline
		E001                      & Error al guardar configuración & Mostrar mensaje: 'No se pudo guardar la configuración'      \\ \hline
		E002                      & Campo obligatorio no definido  & Mostrar advertencia: 'Este campo no puede quedar en blanco' \\ \hline
	\end{tabular}
\end{table}

\newpage

\subsubsection*{Estado}
En definición


\subsubsection*{Requerimientos suplementarios}

\begin{table}[h]
	\centering
	\caption{Requerimientos Suplementarios}
	\begin{tabular}{|c|p{4cm}|p{8cm}|}
		\hline
		\textbf{Código} & \textbf{Requerimiento} & \textbf{Descripción}                                               \\ \hline
		RS01            & Seguridad              & El sistema debe encriptar contraseñas y usar HTTPS                 \\ \hline
		RS02            & Disponibilidad         & El sistema debe estar disponible al menos 95\% del tiempo          \\ \hline
		RS03            & Usabilidad             & La interfaz debe ser intuitiva y accesible                         \\ \hline
		RS04            & Rendimiento            & El cálculo de nómina no debe superar 5 segundos para 100 empleados \\ \hline
		RS05            & Escalabilidad          & El sistema debe permitir agregar nuevos módulos                    \\ \hline
		RS06            & Soporte                & El sistema debe de ofrecer documentación de ayuda integrada        \\ \hline
	\end{tabular}
\end{table}

\subsection*{Requerimientos de documentación}

Esta sección describe la documentación que acompañará al sistema de gestión de nóminas.
Se detallan los manuales, guías de instalación y otros recursos necesarios para garantizar
que los usuarios y administradores del sistema puedan utilizarlo, instalarlo y mantenerlo
adecuadamente. Además, se especifica el formato, contenido y medios de distribución de
dicha documentación.

\subsubsection*{Manual de Usuario}

El sistema contará con un manual de usuario dirigido al personal administrativo y de recursos humanos de \textit{Café y Vivero Verde Pradera}. Este manual incluirá:

\begin{itemize}
	\item Introducción al sistema y sus módulos.
	\item Guía paso a paso para cada funcionalidad (registro de empleados, cálculo de nómina, reportes, etc.).
	\item Imágenes de la interfaz.
	\item Glosario de términos y símbolos utilizados.
	\item Resolución de problemas comunes.
\end{itemize}

El manual estará disponible en formato PDF y como parte de la sección de ayuda en línea. Será conciso, visual y redactado en español.

\subsubsection*{Ayuda en Línea}

El sistema integrará un módulo de ayuda contextual. Esta ayuda se presentará como:

\begin{itemize}
	\item Íconos de ayuda ``?'' cerca de cada formulario o sección.
	\item Ventanas emergentes con descripciones breves.
	\item Acceso a una sección completa de preguntas frecuentes (FAQ).
\end{itemize}

Todo el contenido estará disponible sin conexión a internet y se mantendrá actualizado en cada versión del sistema.

\subsubsection*{Guías de Instalación, Configuración y Archivo \texttt{README}}

Se incluirá un documento de instalación técnica para el equipo de desarrollo o soporte técnico de la empresa. Contendrá:

\begin{itemize}
	\item Requisitos del sistema.
	\item Pasos de instalación del \textit{backend} y \textit{frontend}.
	\item Configuración de la base de datos y variables de entorno.
	\item Archivo \texttt{README.md} con:
	      \begin{itemize}
		      \item Novedades de cada versión.
		      \item Cambios importantes en funcionalidades.
		      \item Compatibilidad hacia versiones anteriores.
	      \end{itemize}
\end{itemize}

\subsubsection*{Etiquetado y Empaquetado}

\begin{itemize}
	\item El sistema incluirá elementos visuales consistentes con la identidad gráfica de la empresa, como logotipo y esquema de colores.
	\item Toda la documentación llevará la marca del sistema y datos de contacto del desarrollador.
	\item Los íconos del sistema seguirán un estilo limpio, con buena accesibilidad visual.
	\item Se mostrará el aviso de derechos reservados en la pantalla de \textit{login}.
\end{itemize}









\section{Cronograma de Desarrollo y Pruebas}

\begin{table}[h]
	\centering
	\caption{Cronograma de Desarrollo y Pruebas del Proyecto}
	\renewcommand{\arraystretch}{1.2}
	\begin{tabular}{|p{4.2cm}|p{2.5cm}|p{2.5cm}|p{2.2cm}|p{3.5cm}|}
		\hline
		\textbf{Sprint / Actividad}             & \textbf{Fecha de Inicio} & \textbf{Fecha de Finalización} & \textbf{Duración} & \textbf{Periodo}                \\ \hline
		Sprint 1: Setup del Proyecto            & 15 mayo 2025             & 31 mayo 2025                   & 2 semanas         & Mayo 2025 (segunda mitad)       \\ \hline
		Sprint 2: Módulo de Login               & 1 junio 2025             & 15 junio 2025                  & 2 semanas         & Junio 2025 (primera mitad)      \\ \hline
		Sprint 3: Módulo de Interfaz de Usuario & 16 junio 2025            & 15 julio 2025                  & 4 semanas         & Junio - Julio 2025              \\ \hline
		Sprint 4: Admin. de Eventos Laborales   & 16 julio 2025            & 15 agosto 2025                 & 4 semanas         & Julio - Agosto 2025             \\ \hline
		Sprint 5: Cálculo de Nóminas            & 16 agosto 2025           & 15 septiembre 2025             & 4 semanas         & Agosto - Septiembre 2025        \\ \hline
		Sprint 6: Generación de Reportes        & 16 septiembre 2025       & 30 septiembre 2025             & 2 semanas         & Septiembre 2025 (segunda mitad) \\ \hline
		Sprint 7: Módulo de Seguridad           & 1 octubre 2025           & 31 octubre 2025                & 4 semanas         & Octubre 2025                    \\ \hline
		Sprint 8: Pruebas y Entrega             & 1 febrero 2026           & 30 junio 2026                  & 5 meses           & Febrero - Junio 2026            \\ \hline
		\rowcolor{gray!10}
		\textbf{Etapa de Pruebas Funcionales}   & 1 febrero 2026           & 30 abril 2026                  & 3 meses           & Febrero - Abril 2026            \\ \hline
		\rowcolor{gray!10}
		\textbf{Revisión y Documentación Final} & 1 mayo 2026              & 30 junio 2026                  & 2 meses           & Mayo - Junio 2026               \\ \hline
	\end{tabular}
\end{table}



